\documentclass[11pt]{article}
\usepackage{../../../tex/lecturestyle}

\title{Min-Plus Convolution of Convex Sequences}
\author{}
\date{}

\begin{document}
\thispagestyle{fancy}

\lecturetitleband{Miscellaneous}{Min-Plus Convolution of Convex Sequences}{4}

\section{Overview}
For sequences $A=(A_0,\ldots,A_N)$ and $B=(B_0,\ldots,B_M)$, their
\emph{min-plus convolution} is $C_k = \min_{i+j=k}(A_i+B_j)$. Computed
directly, this costs $O(NM)$, and doing much better in general is
believed to be hard. But when both $A$ and $B$ are \emph{convex}, the
whole computation collapses to $O(N+M)$ --- and, reinterpreted through a
simple optimization story, the reason why becomes almost obvious.

\section{Convex Sequences}
\begin{definitionbox}
A sequence $A=(A_0,\ldots,A_N)$ is \emph{convex below} (or just
\emph{convex}) if $A_{i+1}-A_i \le A_{i+2}-A_{i+1}$ for every valid $i$
--- its successive differences are nondecreasing. The reverse inequality
makes it \emph{convex above} (or \emph{concave}). Plotting $(i, A_i)$,
convexity below is exactly a nondecreasing slope as you move right.
\end{definitionbox}

\begin{examplebox}
$(0,0,0,1,2,4)$ is convex below (differences $0,0,1,1,2$, nondecreasing).
$(4,3,2,1,0)$ is both convex below and convex above (constant
difference $-1$). $(3,1,3,0,2)$ is neither.
\end{examplebox}

\section{The Main Theorem}
\begin{theorembox}
If $A=(A_0,\ldots,A_N)$ and $B=(B_0,\ldots,B_M)$ are both convex below,
their min-plus convolution $C$ can be computed in $O(N+M)$ time, and $C$
is itself convex below.

\emph{Proof.} Adding a constant to every $A_i$ (or every $B_j$) just adds
that constant to every $C_k$, so assume $A_0=B_0=0$. Let $a_i = A_{i+1} -
A_i$ and $b_j = B_{j+1}-B_j$ be the difference sequences; convexity says
both are nondecreasing.

Reinterpret: put $N$ items of costs $a_0,\ldots,a_{N-1}$ in bag $A$, and
$M$ items of costs $b_0,\ldots,b_{M-1}$ in bag $B$. Because the item
costs in each bag are sorted, the cheapest way to pick exactly $i$ items
from bag $A$ is to take the $i$ cheapest, with total cost exactly $A_i$
(this is where convexity is used --- without sorted costs, ``the cheapest
$i$ items'' and ``$A_i$'' need not coincide). Likewise $B_j$ is the
cheapest cost of $j$ items from bag $B$.

Then $A_i + B_j$ is the cheapest cost of picking $i$ items from $A$ and
$j$ from $B$, and $C_k=\min_{i+j=k}(A_i+B_j)$ is the cheapest cost of
picking $k$ items total from the combined pool of $N+M$ items --- which
is achieved, exactly as in the single-bag case, by taking the $k$
globally cheapest items. So $C_k$ is the $k$-th prefix sum of the
merged, sorted difference sequence $a_0,\ldots,a_{N-1},b_0,\ldots,b_{M-1}$.
Since $a$ and $b$ are each already sorted, merging them costs only
$O(N+M)$ (an ordinary two-pointer merge, as in merge sort), and so does
the resulting prefix sum. Undoing the initial normalization just adds
back $A_0+B_0$.

$C$'s difference sequence is, by construction, this merged (hence
sorted) sequence, so $C$ is itself convex below. \qed
\end{theorembox}

\begin{notebox}
The same argument, run on convex-above sequences with a max-plus
convolution, works identically (pick the $k$ \emph{most expensive} items
instead).
\end{notebox}

% FIGURE: mp_merge
\begin{figure}[H]
\centering
\begin{tikzpicture}[x=1.05cm,y=1.0cm]
\node[llabel,anchor=east] at (-0.7,3.0) {slopes of $a$};
\node[lcell,lhi,minimum width=0.95cm] at (0,3.0) {$-3$};
\node[lcell,lhi,minimum width=0.95cm] at (1,3.0) {$-1$};
\node[lcell,lhi,minimum width=0.95cm] at (2,3.0) {$2$};
\node[llabel,anchor=east] at (-0.7,2.0) {slopes of $b$};
\node[lcell,lhi2,minimum width=0.95cm] at (0,2.0) {$-2$};
\node[lcell,lhi2,minimum width=0.95cm] at (1,2.0) {$0$};
\node[lcell,lhi2,minimum width=0.95cm] at (2,2.0) {$1$};
\node[lcell,lhi2,minimum width=0.95cm] at (3,2.0) {$4$};
\draw[larrow] (7.9,2.5) -- (7.9,1.6) node[midway,right,llabel] {merge into one sorted list};
\node[llabel,anchor=east] at (-0.7,1.1) {merged slopes};
\node[lcell,lhi,minimum width=0.95cm] at (0,1.1) {$-3$};
\node[lcell,lhi2,minimum width=0.95cm] at (1,1.1) {$-2$};
\node[lcell,lhi,minimum width=0.95cm] at (2,1.1) {$-1$};
\node[lcell,lhi2,minimum width=0.95cm] at (3,1.1) {$0$};
\node[lcell,lhi2,minimum width=0.95cm] at (4,1.1) {$1$};
\node[lcell,lhi,minimum width=0.95cm] at (5,1.1) {$2$};
\node[lcell,lhi2,minimum width=0.95cm] at (6,1.1) {$4$};
\draw[larrow] (7.9,0.7) -- (7.9,-0.2) node[midway,right,llabel] {prefix sums, starting from $a_0+b_0=7$};
\node[llabel,anchor=east] at (-0.7,-0.6) {$c=a\oplus b$};
\node[lcell,minimum width=0.95cm,lgray] at (0,-0.6) {$7$};
\node[lcell,minimum width=0.95cm,lgray] at (1,-0.6) {$4$};
\node[lcell,minimum width=0.95cm,lgray] at (2,-0.6) {$2$};
\node[lcell,minimum width=0.95cm,lgray] at (3,-0.6) {$1$};
\node[lcell,minimum width=0.95cm,lgray] at (4,-0.6) {$1$};
\node[lcell,minimum width=0.95cm,lgray] at (5,-0.6) {$2$};
\node[lcell,minimum width=0.95cm,lgray] at (6,-0.6) {$4$};
\node[lcell,minimum width=0.95cm,lgray] at (7,-0.6) {$8$};
\end{tikzpicture}
\caption{Min-plus convolution of convex sequences is a merge. For $a$ with slopes $(-3,-1,2)$ and $b$ with slopes $(-2,0,1,4)$, merging the two sorted slope lists and taking prefix sums from $a_0+b_0=7$ yields $c=(7,4,2,1,1,2,4,8)$, which equals $\min_{i+j=k}(a_i+b_j)$ for every $k$.}\label{fig:mp_merge}
\end{figure}

\section{A DP Speedup Application}
Convex min-plus convolution shows up disguised inside ordinary dynamic
programs, and recognizing it can turn an $O(N^2)$ solution into an
$O(N \log N)$ one.

\begin{definitionbox}
Given $A=(A_1,\ldots,A_N)$, for every $k=0,\ldots,N$ find the minimum
possible sum of exactly $k$ chosen elements of $A$.
\end{definitionbox}

The direct answer is immediate: sort $A$ and take prefix sums, in
$O(N\log N)$. But a natural DP formulation is also tempting: let
$\mathrm{dp}[i]$ be the answer using only the first $i$ processed
elements, updated one element $x$ at a time by
$\mathrm{dp}'[i]=\min(\mathrm{dp}[i], \mathrm{dp}[i-1]+x)$ (skip or take
$x$). Implemented as a growing array, this costs $\Theta(N^2)$ overall.

The connection: each such update is exactly a min-plus convolution of the
current $\mathrm{dp}$ table with the two-element sequence $(0, x)$, which
is trivially convex. Since a convex sequence convolved with a convex
sequence stays convex, $\mathrm{dp}$ is convex at every step --- and by
the theorem, updating it means merging a single value $x$ into
$\mathrm{dp}$'s (sorted) difference sequence. Doing this $N$ times, for
$x = A_1,\ldots,A_N$ in turn, is nothing more than inserting all of $A$'s
values one at a time into a sorted sequence: exactly sorting $A$. The
$O(N^2)$-looking DP and the $O(N\log N)$ greedy algorithm are, under this
lens, the same algorithm.

\begin{notebox}
This particular problem is simple enough that the connection mostly
confirms what greedy intuition already suggested. The real payoff shows
up in less obvious DPs where a min-plus convolution with a convex
sequence is buried inside a larger recurrence (tree DPs merging convex
child tables are a common source) --- recognizing convexity there is what
unlocks a fast solution.
\end{notebox}

\section{Reference Implementations}
Both examples below implement the naive $O(NM)$ definition and the
$O(N+M)$ convex-specific algorithm, verify them against both of the
lecture's own worked examples and against each other on random convex
sequences, confirm the result stays convex, and work through the DP
speedup application above.

\codenote
\subsection{Rust}
\corelisting{Rust}{example.rs}

\subsection{Python}
\corelisting{Python}{example.py}

\section{Worked Example}
\begin{examplebox}
For $A=(2,0,0,1,2)$ and $B=(0,0,1,3)$ (both convex, with difference
sequences $(-2,0,1,1)$ and $(0,1,2)$ respectively), merging the
differences gives $(-2,0,0,1,1,1,2)$, and accumulating from $A_0+B_0=2$
gives $C=(2,0,0,0,1,2,3,5)$ --- matching both the lecture's own stated
result and the naive $O(NM)$ computation exactly, as confirmed by both
reference implementations.
\end{examplebox}

\section{Further Reading}
\begin{itemize}[leftmargin=1.5em]
  \item AtCoder Algorithm Lectures index (source of this curriculum):
        \url{https://info.atcoder.jp/entry/algorithm_lectures/index}
  \item M. Cygan, M. Mucha, K. W\k{e}grzycki, M. W{\l}odarczyk, ``On
        Problems Equivalent to (min,+)-Convolution'', ACM Transactions
        on Algorithms 15(1), 2019, pp.\ 1--25.
\end{itemize}

\end{document}
